\documentclass[11pt]{article}
\usepackage{geometry}
\usepackage{datatool} % For CSV parsing and math
\usepackage{booktabs} % For beautiful table lines

\geometry{a4paper, margin=1in}

% --- Configuration ---
\newcommand{\vatrate}{0.20} % 20% VAT
% We use a separate file for the customer name so our Python script can inject it easily
\newcommand{\CustomerName}{\input{current_customer.txt}} 

\begin{document}

% Header
\noindent
\begin{minipage}{0.5\textwidth}
    {\Huge \textbf{INVOICE}}\\
    \vspace{0.5em}\\
    \textbf{Billed To:}\\
    \CustomerName
\end{minipage}
\begin{minipage}{0.5\textwidth}
    \raggedleft
    \textbf{Date:} \today\\
    \textbf{Invoice \#:} \DTLcurrentpage % Placeholder for invoice number
\end{minipage}

\vspace{3em}

% Load the CSV database
\DTLloaddb{items}{items.csv}

% Initialize our global SubTotal variable to 0
\gdef\SubTotal{0}

% The Invoice Table
\noindent
\begin{tabular*}{\textwidth}{@{\extracolsep{\fill}} l r r r @{}}
\toprule
\textbf{Description} & \textbf{Qty} & \textbf{Unit Price (\$)} & \textbf{Total (\$)} \\
\midrule
% Loop through the CSV
\DTLforeach{items}{\desc=Item,\qty=Quantity,\price=UnitPrice}{%
  % 1. Multiply Qty by Price to get the Row Total
  \DTLmul{\rowtotal}{\qty}{\price}%
  \DTLround{\rowtotal}{\rowtotal}{2}% Round to 2 decimal places
  
  % 2. Add Row Total to the global SubTotal
  \DTLgadd{\SubTotal}{\rowtotal}%
  
  % 3. Print the row in the table
  \desc & \qty & \price & \rowtotal \\
}
\bottomrule
\end{tabular*}

\vspace{2em}

% --- Final Calculations ---
% Calculate VAT amount
\DTLmul{\VatAmount}{\SubTotal}{\vatrate}
\DTLround{\VatAmount}{\VatAmount}{2}

% Calculate Grand Total
\DTLadd{\GrandTotal}{\SubTotal}{\VatAmount}
\DTLround{\GrandTotal}{\GrandTotal}{2}

% Summary Box
\begin{flushright}
\begin{tabular}{r l}
\textbf{Subtotal:} & \$\SubTotal \\
\textbf{VAT (20\%):} & \$\VatAmount \\
\midrule
\textbf{Grand Total:} & \textbf{\$\GrandTotal} \\
\end{tabular}
\end{flushright}

\end{document}